\documentclass[12pt, a4paper]{article}

% --- Packages ---
\usepackage[utf8]{inputenc}
\usepackage[T1]{fontenc}
\usepackage{amsmath, amssymb, amsthm}
\usepackage{graphicx}
\usepackage{hyperref}
\usepackage{natbib}
\usepackage{booktabs}
\usepackage{caption}
\usepackage{subcaption}
\usepackage{geometry}
\usepackage{xcolor}
\usepackage{algorithm}
\usepackage{algorithmic}

\geometry{margin=1in}

% --- Title ---
\title{Topological Signatures of Microbiome Dysbiosis and Their Predictive \\
  Relationship to Gut-Brain Axis Dysfunction: \\
  A Persistent Homology Approach}

\author{Adam [Last Name] \\
  \textit{Institution} \\
  \texttt{email@institution.edu}}

\date{\today}

% --- Document ---
\begin{document}

\maketitle

\begin{abstract}
The human gut microbiome exhibits complex ecological dynamics whose disruption---dysbiosis---drives a pathological cascade from intestinal barrier degradation through systemic immune activation to neuroinflammation and gut-brain axis dysfunction. While traditional analyses focus on compositional shifts in microbial taxa, the higher-order structural organization of microbial communities remains poorly characterized. Here, we apply persistent homology from topological data analysis (TDA) to co-occurrence networks derived from 16S rRNA sequencing data from the Human Microbiome Project and American Gut Project. We identify topological signatures---persistent loops and voids in the co-occurrence landscape---that distinguish healthy microbiome configurations from dysbiotic states. Our sliding-window analysis reveals abrupt topological regime transitions that precede compositional shifts, paralleling findings from topological analysis of financial market instability. These topological biomarkers correlate with intestinal permeability markers (zonulin, LPS-binding protein), short-chain fatty acid levels (butyrate), and neuropeptide concentrations (serotonin, VIP, CCK), providing a novel structural framework linking microbial network topology to the dysbiosis-permeability-neuroinflammation cascade. We further examine mycotoxin-induced barrier dysfunction as a case study of environmentally triggered topological regime shifts, with implications for understanding chronic inflammatory response syndrome (CIRS) and mold-associated illness.
\end{abstract}

\section{Introduction}
\label{sec:introduction}

% (~4-5 pages)

% Opening: The gut microbiome as a complex adaptive system
% - 100 trillion microorganisms, 1000+ species
% - Not just passengers — active participants in host physiology
% - Metabolic, immune, and neurological functions

% The dysbiosis problem
% - Definition: disruption of the symbiotic equilibrium
% - Clinical relevance: IBD, obesity, diabetes, neurological disorders
% - The cascade: dysbiosis → SCFA depletion → barrier degradation →
%   endotoxin translocation → systemic inflammation → neuroinflammation

% Limitations of current approaches
% - Alpha/beta diversity metrics capture "how much" but not "what structure"
% - Differential abundance (DESeq2) identifies individual taxa but misses
%   community-level organization
% - Co-occurrence networks capture pairwise relationships but not
%   higher-order topology

% The TDA opportunity
% - Persistent homology captures multi-scale topological features
% - Proven in other complex systems: financial markets, protein structure,
%   neural networks, materials science
% - Key parallel: correlation homogeneity → topological stability
%   (markets ↔ microbial ecosystems)

% The gut-brain axis connection
% - Vagal nerve as bidirectional communication highway
% - Neuropeptides produced/modulated by gut microbiota:
%   serotonin (95% produced in gut), ghrelin, CCK, VIP
% - Intestinal permeability as the gateway mechanism
% - Environmental triggers: mycotoxins, antibiotics, diet

% Our contribution
% - First application of persistent homology to detect regime shifts
%   in gut microbial co-occurrence networks
% - Topological biomarkers that correlate with permeability markers
%   and neuropeptide levels
% - Sliding-window regime detection as early warning system
% - Mycotoxin-induced dysbiosis as case study (CIRS connection)

% Paper outline
% Section 2: Background on microbiome, TDA, gut-brain axis
% Section 3: Theoretical framework linking topology to ecology
% Section 4: Methods (datasets, pipeline, statistics)
% Section 5: Results
% Section 6: Discussion (interpretation, CIRS implications, limitations)
% Section 7: Conclusion

\section{Background}
\label{sec:background}

% (~8-10 pages)

\subsection{The Human Gut Microbiome}
% Microbial ecology of the gut
% Key phyla: Firmicutes, Bacteroidetes, Proteobacteria, Actinobacteria
% Functional roles: nutrient metabolism, immune training, pathogen resistance
% The Firmicutes/Bacteroidetes ratio as a crude health indicator
% Why composition alone is insufficient — need structural understanding

\subsection{Microbiome Dysbiosis and the Permeability Cascade}
\label{subsec:permeability_cascade}
% Definition: loss of beneficial organisms, expansion of pathobionts,
%   reduction in overall diversity
% The cascade mechanism:
%   1. Dysbiosis → reduced butyrate-producing bacteria (Faecalibacterium,
%      Roseburia, Eubacterium)
%   2. Butyrate depletion → reduced tight junction protein expression
%      (claudins, occludin, ZO-1)
%   3. Barrier degradation → increased intestinal permeability ("leaky gut")
%   4. LPS and other endotoxins translocate to systemic circulation
%   5. TLR4 activation → NF-κB → pro-inflammatory cytokines
%      (TNF-α, IL-1β, IL-6)
%   6. Systemic inflammation → neuroinflammation via blood-brain barrier
%      and vagal afferents
% Permeability biomarkers: zonulin, LPS-binding protein (LBP),
%   intestinal fatty acid-binding protein (I-FABP),
%   lactulose/mannitol ratio
% Associated diseases: IBD, IBS, obesity, T2D, depression, anxiety,
%   Parkinson's, Alzheimer's

\subsection{The Gut-Brain Axis}
\label{subsec:gut_brain}
% Bidirectional communication via:
%   - Vagus nerve (primary conduit)
%   - Enteric nervous system (ENS, "second brain")
%   - Immune signaling (cytokines crossing BBB)
%   - Microbial metabolites (SCFAs, tryptophan metabolites)
% Neuropeptides modulated by gut microbiota:
%   - Serotonin (5-HT): 95\% produced by enterochromaffin cells in gut,
%     microbiota regulate tryptophan hydroxylase expression
%   - Ghrelin: appetite regulation, modulated by gut composition
%   - Cholecystokinin (CCK): satiety, gut motility
%   - Vasoactive intestinal peptide (VIP): anti-inflammatory,
%     neuroprotective, barrier maintenance
%   - GABA: produced by Lactobacillus and Bifidobacterium spp.
% Disrupted signaling in dysbiosis → cognitive impairment,
%   emotional dysregulation, altered decision-making

\subsection{Mycotoxins and Environmental Triggers of Dysbiosis}
\label{subsec:mycotoxins}
% Mycotoxins as a concrete environmental trigger
% Trichothecenes (T-2 toxin, deoxynivalenol/DON):
%   - Directly reduce tight junction protein expression
%   - Decrease mucin-secreting goblet cells
%   - Alter microbiome composition (reduce Lactobacillus, increase
%     Proteobacteria)
% Aflatoxins, ochratoxin A: additional gut barrier effects
% Chronic Inflammatory Response Syndrome (CIRS):
%   - Biotoxin illness from water-damaged buildings
%   - Involves gut barrier dysfunction, immune dysregulation,
%     neuroinflammation
%   - HLA-DR genetic susceptibility (~24\% of population)
%   - Shoemaker protocol: cholestyramine, VIP nasal spray
% Mycotoxin-induced dysbiosis → bacterial translocation →
%   chronic inflammation → neurological effects
% This provides a concrete use case for topological regime detection

\subsection{Co-occurrence Networks in Microbial Ecology}
\label{subsec:cooccurrence}
% Why networks: taxa don't exist in isolation
% Network construction from abundance data:
%   - Pearson/Spearman correlation (problematic for compositional data)
%   - SparCC: handles compositionality via log-ratio variance
%   - SPIEC-EASI: sparse inverse covariance estimation
% Network metrics: modularity, clustering coefficient, hub scores,
%   betweenness centrality
% Ecological interpretation: modules = functional guilds,
%   hubs = keystone species
% Limitations: pairwise relationships miss higher-order structure

\subsection{Topological Data Analysis}
\label{subsec:tda}
% Simplicial complexes: generalizing graphs to higher dimensions
%   - 0-simplices (vertices), 1-simplices (edges),
%     2-simplices (triangles), etc.
% Filtrations: growing simplicial complexes by a parameter
%   - Vietoris-Rips: add simplex when all pairwise distances < ε
%   - \v{C}ech: add simplex when balls of radius ε have common intersection
% Persistent homology: tracking topological features across filtration
%   - Birth and death of features
%   - Persistence = death - birth (long-lived features are "real")
% Betti numbers and their ecological interpretation:
%   - β₀: connected components (isolated community clusters)
%   - β₁: loops (ecological feedback cycles, mutualism rings)
%   - β₂: voids (higher-order community cavities)
% Persistence diagrams, barcodes, landscapes
% Prior applications:
%   - Protein structure and function \citep{camara2017topological}
%   - Neural spike train analysis
%   - Financial market regime detection (the parallel to this work)
%   - Biological aggregation models \citep{topaz2015topological}
%   - Zebrafish pattern formation \citep{mcguirl2020topological}

\section{Theoretical Framework}
\label{sec:framework}

\subsection{From Abundance Tables to Topological Spaces}
% OTU table → correlation/distance matrix → simplicial complex
% Choice of filtration and its biological meaning

\subsection{Topological Invariants as Ecological Descriptors}
% H0: connected components (community clusters)
% H1: loops (ecological cycles, feedback loops)
% H2: voids (higher-order community organization)

\subsection{Regime Detection via Sliding Window Topology}
% Temporal/gradient sliding windows
% Tracking topological features over parameter changes
% Defining topological regime transitions

\section{Methods}
\label{sec:methods}

\subsection{Datasets}
% HMP Phase 1: sample selection, preprocessing
% AGP: sample selection, metadata variables

\subsection{Data Preprocessing}
% Quality filtering: prevalence and read depth thresholds
% CLR transformation for compositionality
% Genus-level aggregation

\subsection{Co-occurrence Network Construction}
% SparCC correlation estimation
% Thresholding and edge weighting
% Network validation

\subsection{Persistent Homology Pipeline}
% Vietoris-Rips filtration on correlation distance
% Ripser computation parameters
% Feature extraction: Betti curves, landscapes, entropy

\subsection{Statistical Analysis}
% Permutation tests for group comparisons
% Correlation with biomarkers
% Machine learning classification

\section{Results}
\label{sec:results}

% (~8-10 pages)

\subsection{Dataset Characteristics and Preprocessing}
% Sample counts after QC for HMP, AGP, and curatedMetagenomicData
% NT pathway annotation: genera mapped per pathway
% NT production score distributions across cohorts
% Table 1: Dataset summary statistics

\subsection{Co-occurrence Network Structure and NT Subnetworks}
\label{subsec:network_results}
% Full network topology across health states
% NT-specific subnetwork comparisons:
%   - Serotonin subnetwork: Enterococcus-Clostridium-Ruminococcus structure
%   - GABA subnetwork: Lactobacillus-Bifidobacterium-Bacteroides connectivity
%   - Dopamine subnetwork: smaller, Enterococcus-Escherichia core
%   - SCFA subnetwork: Faecalibacterium-Roseburia-Eubacterium hub structure
% Figure: NT subnetwork visualizations

\subsection{Persistence Diagrams Reveal Topological Signatures of NT Production}
\label{subsec:persistence_results}
% Full network: H$_0$, H$_1$, H$_2$ features
% NT subnetworks: pathway-specific topological profiles
%   - High-serotonin samples: more persistent H$_1$ loops in serotonin subnetwork
%   - Cross-feeding loops between SCFA and serotonin subnetworks
% Betti curves and persistence landscapes as NT-state discriminators
% Figure: Persistence diagrams, Betti curves for high vs. low NT groups

\subsection{Mycobiome Disruption Fragments NT-Network Topology}
\label{subsec:mycobiome_results}
% High-fungal vs. low-fungal burden comparison:
%   - Loss of H$_1$ features in GABA subnetwork (Candida displaces Lactobacillus)
%   - Fragmentation of H$_0$ in serotonin subnetwork
%   - Disruption scores correlate with topological degradation
% Per-fungal-genus effects: Candida vs. Aspergillus vs. Malassezia
% Figure: Topological comparison high vs. low fungal burden

\subsection{AMR Carriers as Topological Disruptors}
\label{subsec:amr_results}
% AMR burden correlates with NT-network topology metrics
% Resilience simulation: topological degradation curve
%   - NT cross-feeding loops collapse at ~30-40\% non-AMR removal
%   - ``Topological tipping point'' for each NT pathway
% HGT edges: high-HGT AMR carriers form a connected subcomponent
%   that does not participate in NT cross-feeding
% Figure: Resilience curves showing topological degradation

\subsection{Topological Regime Transitions in NT Networks}
\label{subsec:regime_results}
% Sliding window analysis ordered by NT production gradient
% Topological transitions precede compositional shifts
% Figure: Multi-panel regime detection with topological features

\subsection{Topological Features Predict NT Production and Mood Outcomes}
\label{subsec:prediction_results}
% Correlation with NT production scores
% Correlation with AGP self-reported mental health outcomes
% Cross-cohort validation with curatedMetagenomicData
% ML classification: topological features vs. diversity metrics
% Table: Classification performance comparison

\section{Discussion}
\label{sec:discussion}

% Interpretation of topological signatures
% Biological meaning of persistent loops in co-occurrence networks
% Comparison with existing dysbiosis indicators
% Early warning signals from topological transitions
% Implications for the gut-brain axis
% Limitations and caveats
% Future directions

\section{Conclusion}
\label{sec:conclusion}

% (~1-2 pages)

% Summary of key findings:
%   1. Persistent homology reveals distinct topological signatures in
%      microbial co-occurrence networks associated with neurotransmitter
%      production capacity (serotonin, GABA, dopamine, SCFAs)
%   2. NT-producing community topology — particularly persistent H$_1$ loops
%      representing cross-feeding cycles — predicts NT production potential
%      better than taxonomic composition alone
%   3. Mycobiome (fungal) co-colonization fragments NT-network topology,
%      especially the GABA and serotonin pathways, providing a topological
%      mechanism for mold-associated neurological symptoms
%   4. AMR carriers act as topological disruptors with a quantifiable
%      ``tipping point'' beyond which NT cross-feeding networks collapse
%   5. Topological regime transitions precede compositional shifts in
%      NT-associated taxa, supporting topology as an early warning system

% Contributions:
%   - First application of persistent homology (not Mapper) to microbiome
%     co-occurrence networks with Betti curves and persistence landscapes
%   - First linking of topological features to a functional outcome
%     (neurotransmitter production)
%   - First topological characterization of mycobiome-bacterial and
%     AMR-bacterial disruption
%   - Bridges TDA-microbiome and TDA-neuroscience literatures

% Broader implications:
%   - Gut-brain axis: community topology, not just composition,
%     determines NT signaling capacity
%   - Clinical: topological monitoring during antibiotic treatment,
%     probiotic interventions, environmental exposure management
%   - Personalized medicine: topological profiles could guide
%     targeted interventions to restore NT-producing network structure

% Closing:
%   The gut microbiome is not merely a collection of NT-producing species ---
%   it is a structured ecological network whose topology determines its
%   collective capacity for neurotransmitter biosynthesis. By revealing
%   how this topology is disrupted by fungi, antibiotics, and resistant
%   organisms, and how it correlates with neurotransmitter output and
%   neurological outcomes, we open a new structural avenue for understanding
%   and optimizing the gut-brain axis.


\bibliographystyle{plainnat}
\bibliography{references}

\end{document}
